\section{Related Work}
\label{sec:related}

\paragraph{End-to-end document understanding.}
Generative document understanding has converged on encoder--decoder
transformers that consume the page image and emit structured output.
DONUT~\cite{kim2022donut} couples a Swin~\cite{liu2021swin} encoder
with a BART~\cite{lewis2020bart} decoder; Pix2Struct~\cite{lee2023pix2struct}
and Dessurt~\cite{davis2022dessurt} take similar approaches.  Each
trades transparency of the field-assignment step for end-to-end
optimisation.  Our DONUT arm follows the original SROIE recipe
exactly~\cite{kim2022donut}.

\paragraph{Modular detect-then-read pipelines.}
The classic detect-then-read paradigm separates text-line detection
(YOLOv8~\cite{jocher2023yolov8}, DBNet~\cite{liao2020dbnet},
Tesseract~\cite{smith2007tesseract}) from text recognition
(TrOCR~\cite{li2023trocr}, PaddleOCR~\cite{paddleocr}).  Field
assignment historically falls to either a learned head
(PICK~\cite{yu2020pick}) or a deterministic regex-and-spatial-rules
component.  We evaluate the latter as the engineering reference.

\paragraph{Receipt and form benchmarks.}
SROIE~\cite{huang2019icdar} (626 receipts, 4 fields) is the standard
single-vendor benchmark.  CORD~\cite{park2019cord} extends to Korean
restaurant receipts with a richer schema; WildReceipt~\cite{sun2021wildreceipt}
extends to in-the-wild imaging.  Form benchmarks adjacent to receipts
include FUNSD~\cite{jaume2019funsd}, RVL-CDIP~\cite{harley2015rvlcdip},
DocVQA~\cite{mathew2021docvqa}, and Kleister~\cite{stanislawek2021kleister}.
We evaluate on canonical SROIE Task-3 only.

\paragraph{Replication and methodology in document AI.}
A growing literature documents the gap between published numbers
and reproducible numbers in document AI.  Our 14-bug catalogue
(Sec.~\ref{sec:bugs}) sits in this lineage as a per-bug-mechanism
record of what silently breaks during reproduction of a published
recipe.  The reproducibility methodology mirrors NeurIPS-2024
checklist conventions.

\paragraph{Carbon and energy reporting.}
We report per-stage GPU energy (kWh) and CO\textsubscript{2}eq (kg)
following Strubell et al.~\cite{strubell2019energy} as a routine
methodological practice; emission accounting is not a contribution
of this paper.
